\section{Proofs}
\label{app:proofs}

\subsection{Exactness of the residual scheme}
\label{app:exactness}

We record the identity the whole measurement apparatus is anchored on. Let $q, p$ be densities on
$\R^\dd$, let $\alpha(x) = \min\{1, p(x)/q(x)\}$, let $\bar{A} = \E_{q}[\alpha]$, and let
$Z = \int (p(u) - q(u))_+ \, du$. Draw $x \sim q$, accept with probability $\alpha(x)$, and on
rejection draw from $r(x) = (p(x) - q(x))_+ / Z$. The emitted density is
\[
\tilde{p}(x) = q(x)\alpha(x) + (1 - \bar{A})\,r(x).
\]

First, $Z = 1 - \bar{A}$. Writing $\min\{a,b\} = b - (b-a)_+$ with $a = p$, $b = q$,
\begin{align*}
\bar{A} &= \int q(x)\min\{1,\, p(x)/q(x)\} \, dx
         = \int \min\{q(x), p(x)\} \, dx \\
        &= \int \bigl(q(x) - (q(x) - p(x))_+\bigr) dx
         = 1 - \int (q - p)_+ .
\end{align*}
Since $\int p = \int q = 1$ we have $\int (p-q)_+ = \int (q-p)_+$, so $Z = 1 - \bar A$ as claimed.
Substituting,
\[
\tilde{p}(x) = \min\{q(x), p(x)\} + (p(x) - q(x))_+ = p(x),
\]
because $\min\{a,b\} + (a - b)_+ = a$ pointwise. Hence exact verification with a correct residual
sampler emits the target \emph{exactly}, and any measurement of its fidelity must return zero up to
Monte Carlo error.

This is why we implement the residual by rejection, proposing $y \sim p$ and accepting with
probability $(1 - q(y)/p(y))_+$, which targets $(p-q)_+/Z$, rather than by resampling from $p$.
Resampling from $p$ gives $\tilde p = q\alpha + (1-\bar A)p \ne p$, so exact verification would score
as \emph{lossy} and every relative fidelity number in \S\ref{sec:conservation} would be measuring the
harness rather than the rule. We report the residual sampler's convergence rate alongside every
fidelity measurement; it is $1.000$ throughout the experiments in this paper.

\subsection{Proof of Proposition~\ref{prop:rowspace} (row-space invariance)}

Let $A \in \R^{\kk\times\dd}$ and $Q \in \R^{\kk\times\kk}$ invertible. For a random vector $x$ with
density $f$, the pushforward of $f$ under $x \mapsto QAx$ relates to that under $x \mapsto Ax$ by the
change-of-variables factor $|\det Q|^{-1}$: writing $f_A$ for the density of $Ax$,
\[
f_{QA}(u) = |\det Q|^{-1} f_A(Q^{-1}u).
\]
Evaluating the acceptance ratio for the statistic $QAx$ at a point $x$, with $s = Ax$ and $u = Qs$,
\[
\frac{p_{QA}(Qs)}{q_{QA}(Qs)}
= \frac{|\det Q|^{-1} p_A(s)}{|\det Q|^{-1} q_A(s)}
= \frac{p_A(s)}{q_A(s)} .
\]
The Jacobian cancels, so $\min\{1, p_{QA}/q_{QA}\} = \min\{1, p_A/q_A\}$ at every $x$. $\square$

Two consequences are used in the main text. Any two matrices with the same row space induce identical
acceptance behaviour, so an acceptance rule of this family is a choice of subspace; and each rule may
therefore be reduced to an orthonormal basis without changing anything, which is what makes the
comparison in Table~\ref{tab:collapse} like for like. When $\kk = \dd$ and $A$ is invertible the same
argument gives $p_A/q_A = p/q$, so full-rank projected verification reproduces exact verification
identically, which is the correctness test we apply to every projection implementation.

\subsection{Proof of Proposition~\ref{prop:tv} (acceptance is one minus a total variation)}

Let $s = T(x)$ with pushforwards $p_T, q_T$. The acceptance probability $\min\{1, p_T(s)/q_T(s)\}$
depends on $x$ only through $s$, and $s \sim q_T$ when $x \sim q$, so
\[
\bar{A}_T = \E_{x \sim q}\!\left[\min\!\left\{1, \tfrac{p_T(T(x))}{q_T(T(x))}\right\}\right]
= \int q_T(s) \min\!\left\{1, \tfrac{p_T(s)}{q_T(s)}\right\} ds
= \int \min\{p_T(s),\, q_T(s)\}\, ds .
\]
For any two densities $\int \min\{f, g\} = 1 - \tfrac{1}{2}\int |f - g|$, so
$\bar{A}_T = 1 - \TV(p_T, q_T)$. Taking $T = \mathrm{id}$ recovers the exact rule's
$\bar{A} = 1 - \TV(p,q)$.

For the inequality, write total variation in its variational form,
$\TV(p,q) = \sup_B |P(B) - Q(B)|$ over measurable $B \subseteq \R^\dd$. Every event of the
pushforward is the image of one such $B$: for measurable $C$, $P_T(C) - Q_T(C) = P(T^{-1}C) -
Q(T^{-1}C)$. The supremum defining $\TV(p_T, q_T)$ is therefore taken over the sub-collection
$\{T^{-1}C\} \subseteq \sigma(x)$, giving $\TV(p_T, q_T) \le \TV(p, q)$ and hence
$\bar{A}_T \ge \bar{A}_{\mathrm{exact}}$. $\square$

\paragraph{When projection buys nothing.} The unconstrained supremum is attained at $B^\star = \{p >
q\}$, and it is attained uniquely up to sets on which $p = q$. Equality in the data-processing step
therefore holds exactly when some $T^{-1}C$ agrees with $B^\star$ off $\{p = q\}$, which is the
condition stated in the proposition. A sufficient condition is that $p/q$ be $\sigma(T)$-measurable,
i.e.\ that $T$ be sufficient for the pair $\{p, q\}$: then the exact test is already a function of
$s$ and nothing has been relaxed. Every other statistic strictly raises acceptance, which is why
acceptance on its own cannot rank rules.

\subsection{Proof of Proposition~\ref{prop:split} (divergence is split, never reduced)}

Let $R = \left[\begin{smallmatrix} A \\ A_\perp\end{smallmatrix}\right]$ be orthogonal and write
$(s, u) = (Ax, A_\perp x)$. The map $x \mapsto (s,u)$ is a bijection with unit Jacobian, and KL is
invariant under bijective reparameterisation, so $\KL(p \| q) = \KL(p_{s,u} \| q_{s,u})$. Factor each
joint density as $p_{s,u} = p_A(s)\, p(u \mid s)$ and likewise for $q$. Then
\begin{align*}
\KL(p_{s,u} \| q_{s,u})
&= \int p_A(s) p(u\mid s) \log \frac{p_A(s)\,p(u\mid s)}{q_A(s)\,q(u\mid s)} \, du \, ds \\
&= \int p_A(s) \log\frac{p_A(s)}{q_A(s)} \, ds
 + \int p_A(s) \int p(u \mid s) \log \frac{p(u\mid s)}{q(u \mid s)} \, du \, ds \\
&= \KL(p_A \| q_A) + \E_{p_A}\!\left[\KL\bigl(p(u\mid s) \| q(u \mid s)\bigr)\right],
\end{align*}
using $\int p(u\mid s)\,du = 1$ in the first term. The left-hand side is $\KL(p\|q)$, which does not
mention $A$.

It remains to identify the second term with the emitted law $\tilde{p}_A(s,u) = p_A(s)\,q(u\mid s)$
of Eq.~\ref{eq:emitter}. Because $\tilde{p}_A$ shares the target's verified marginal, the $p_A$
factors cancel inside the logarithm:
\[
\KL(p \,\|\, \tilde{p}_A)
= \int p_A(s)\,p(u \mid s)\, \log \frac{p_A(s)\,p(u\mid s)}{p_A(s)\,q(u \mid s)}\, du\, ds
= \E_{p_A}\!\left[\KL\bigl(p(u\mid s) \,\|\, q(u \mid s)\bigr)\right],
\]
which is the conditional term exactly. $\square$

The identification is what turns Eq.~\ref{eq:split} from an accounting identity into a statement
about fidelity. The quantity that escapes the test is not a residual defined by subtraction; it is
the KL divergence from the target of a distribution one can actually sample, and it is the price of
sampling from that distribution instead of $p$.

Two things are worth drawing out. Nothing here assumes Gaussianity, shared covariance, or any
structure in the disagreement: the identity holds for arbitrary densities and arbitrary orthogonal
splits, which is why we state it before specialising. And both terms are non-negative, so
$\KL(p_A \| q_A) \le \KL(p \| q)$ for every $A$: a projected test can only ever see a fraction of the
divergence, and the remainder is not destroyed but relocated into the conditional term, which is
exactly the fidelity the relaxed rule leaks. What a designer controls is the partition, not the
total.

\subsection{The idealised projected emitter and the one we measure}
\label{app:emitter}

Eq.~\ref{eq:emitter} defines $\tilde{p}_A(s,u) = p_A(s)\,q(u\mid s)$: verify $s$, repair rejections
\emph{within the verified marginal}, and take the complement from the draft. Our experiments do not
implement that sampler. They keep the rejection path of ordinary speculative sampling --- on
rejection, draw from the full-space residual $(p - q)_+$ --- and vary only the acceptance test. Two
reasons. The comparison across the five rules is then a comparison of acceptance tests and nothing
else, and the exact rule ($\kk = \dd$) degenerates to standard speculative sampling, which emits $p$
exactly and so anchors the fidelity axis at a true zero.

The two schemes differ, and the difference is not negligible, so it should be stated rather than
elided. The full-space residual corrects a rejected position in \emph{every} coordinate, including
the unverified ones; $\tilde{p}_A$ corrects only the verified block. Our repair therefore uses
strictly more information about the target and should leak less. Table~\ref{tab:emitter} confirms
this on a shared-covariance Gaussian pair, where every predicted quantity is available in closed
form ($\sigma = 1$, per-coordinate gap $0.578$, $4\times10^5$ samples per cell,
\texttt{scripts/check\_projected\_emitter.py}).

\begin{table}[h]
\centering
\small
\caption{The idealised emitter of Eq.~\ref{eq:emitter} against the scheme our experiments run.
Acceptance matches Proposition~\ref{prop:conservation} in every cell. The idealised emitter's excess
NLL reproduces the conservation identity's leakage term $\|A_\perp\delta\|^2/2\sigma^2$ to three
decimals, which is the numerical check that Eq.~\ref{eq:split} describes a real sampler. Our scheme
leaks strictly less, by a margin that grows with the unverified dimension.}
\label{tab:emitter}
\begin{tabular}{rr rr rrr}
\toprule
& & \multicolumn{2}{c}{acceptance} & \multicolumn{3}{c}{excess NLL (nats)} \\
\cmidrule(lr){3-4}\cmidrule(lr){5-7}
$\dd$ & $\kk$ & measured & predicted & idealised $\tilde{p}_A$ & $\|A_\perp\delta\|^2/2\sigma^2$ & ours \\
\midrule
 16 &  1 & 0.772 & 0.773 &  2.517 &  2.506 & 1.856 \\
 16 &  4 & 0.564 & 0.563 &  2.012 &  2.005 & 1.068 \\
 16 & 16 & 0.248 & 0.248 &  0.005 &  0.000 & 0.001 \\
\midrule
 64 &  1 & 0.772 & 0.773 & 10.521 & 10.524 & 8.064 \\
 64 &  4 & 0.564 & 0.563 & 10.023 & 10.023 & 5.630 \\
 64 & 16 & 0.248 & 0.248 &  8.007 &  8.018 & 2.041 \\
\midrule
128 &  1 & 0.772 & 0.773 & 21.222 & 21.214 & 16.348 \\
128 &  4 & 0.564 & 0.563 & 20.731 & 20.713 & 11.690 \\
128 & 16 & 0.247 & 0.248 & 18.694 & 18.709 & 4.701 \\
\bottomrule
\end{tabular}
\end{table}

Three readings. First, the middle columns agree to three decimals, which is a direct numerical
confirmation that the conditional term of Eq.~\ref{eq:split} is the divergence of an emitter and not
a bookkeeping remainder. Second, at $\kk = \dd$ both schemes sit at zero, as they must: the test is
exact and $\tilde{p}_A = p$. Third, our measured costs are lower everywhere else, substantially so
when much is left unverified --- $4.7$ nats against $18.7$ at $\dd{=}128, \kk{=}16$. The consequence
for the paper's argument is one-directional. Every acceptance--fidelity frontier we report is traced
under the \emph{more} favourable repair, so the obstruction those frontiers exhibit is not an
artefact of a weak sampler; adopting the idealised scheme would move every rule to a higher fidelity
cost at unchanged acceptance.

\subsection{Proof of Proposition~\ref{prop:conservation} (Gaussian conservation)}

Let $p = \mathcal{N}(\mu_p, \sigma^2 I_\dd)$, $q = \mathcal{N}(\mu_q, \sigma^2 I_\dd)$,
$\delta = \mu_p - \mu_q$. Let $A \in \R^{\kk\times\dd}$ have orthonormal rows, so $AA^\top = I_\kk$,
and let $A_\perp \in \R^{(\dd-\kk)\times\dd}$ be an orthonormal basis of the orthogonal complement of
the row space, so that the stacked matrix
$U = \left[\begin{smallmatrix} A \\ A_\perp\end{smallmatrix}\right] \in \R^{\dd\times\dd}$ is
orthogonal.

The pushforward of a Gaussian under a linear map is Gaussian, so
$p_A = \mathcal{N}(A\mu_p,\, A(\sigma^2 I)A^\top) = \mathcal{N}(A\mu_p, \sigma^2 I_\kk)$, and likewise
for $q_A$. Two Gaussians with common covariance $\sigma^2 I_\kk$ and mean gap $A\delta$ satisfy
\[
\KL(p_A \| q_A) = \frac{\|A\delta\|^2}{2\sigma^2}, \qquad
\KL(p_{A_\perp} \| q_{A_\perp}) = \frac{\|A_\perp\delta\|^2}{2\sigma^2}.
\]
Because $U$ is orthogonal it preserves norms, giving
$\|A\delta\|^2 + \|A_\perp\delta\|^2 = \|\delta\|^2$ and therefore
\[
\KL(p_A \| q_A) + \KL(p_{A_\perp} \| q_{A_\perp}) = \frac{\|\delta\|^2}{2\sigma^2} = \KL(p \| q),
\]
independent of $A$.

For the acceptance rate, two Gaussians in $\R^\kk$ with shared covariance $\sigma^2 I$ and mean
separation $m = \|A\delta\|$ have total variation distance $\TV = 2\Phi(m/2\sigma) - 1$, so the
expected acceptance is
\[
\bar{A}_A = 1 - \TV(p_A, q_A) = 2 - 2\Phi\!\left(\frac{\|A\delta\|}{2\sigma}\right),
\]
a strictly decreasing function of $\|A\delta\|^2$ alone. $\square$

\paragraph{Reading the result.} Both the acceptance rate and the divergence that leaks past the test
are determined by the single scalar $\|A\delta\|^2 \in [0, \|\delta\|^2]$: as it increases, acceptance
falls and the leak shrinks, and their sum is pinned. A subspace rule chooses a value of that scalar
and nothing else, which is the precise sense in which the two are traded rather than independently
set, and it is why the adversarial anti-control of \S\ref{sec:conservation} is not a bad rule but a
rule that has chosen a different point on the same curve.

\paragraph{Where the assumptions bind.} Shared covariance is what makes the pushforward covariance
$A\Sigma A^\top$ equal to $\sigma^2 I_\kk$ for \emph{every} orthonormal $A$; with heteroscedastic
components the pushforward acquires a genuine $\kk\times\kk$ covariance that depends on $A$, and the
identity becomes approximate. The mixture structure adds a second source of slack, since the KL
between mixtures does not decompose. Both are reasons to measure rather than assume, which is what
\S\ref{sec:conservation} does; the observed departures from a single curve are $4.0\%$ of the range at
$\dd{=}64$.

\subsection{Proof of Proposition~\ref{prop:random} (random subspaces at finite $\dd$)}
\label{app:random}

An earlier version of this argument took $\E\|A\delta\|^2 = (\kk/\dd)\|\delta\|^2$ and substituted it
into the acceptance formula. That step is not valid. Acceptance is
$f(z) = 2 - 2\Phi(\sqrt{z}/2\sigma)$, which is strictly convex over the relevant range, so
$\E[f(Z)] \neq f(\E Z)$; evaluated over our sweep the substitution understates acceptance by $2.8\%$
to $8.8\%$ (median $5.4\%$). The correct treatment uses the exact law of the captured fraction, and
gives a stronger conclusion than the invalid one did.

Let $A$ be uniform on the Stiefel manifold $V_\kk(\R^\dd)$. By rotational invariance we may instead
fix $A$ and take $\delta$ uniform on the sphere of radius $\|\delta\|$, so
$Z = \|A\delta\|^2/\|\delta\|^2$ is the squared norm of the first $\kk$ coordinates of a uniform unit
vector. Writing that vector as $g/\|g\|$ for $g \sim \mathcal{N}(0, I_\dd)$,
\[
Z = \frac{\sum_{i\le \kk} g_i^2}{\sum_{i \le \dd} g_i^2}
  = \frac{\chi^2_\kk}{\chi^2_\kk + \chi^2_{\dd-\kk}}
  \sim \operatorname{Beta}\!\left(\frac{\kk}{2}, \frac{\dd-\kk}{2}\right),
\]
with the two chi-squared variables independent. Hence, with $M = \|\delta\|/\sigma$,
\[
\E_A\bigl[\bar A_A\bigr] = \E_Z\!\left[2 - 2\Phi\!\left(\frac{M\sqrt{Z}}{2}\right)\right].
\]
Now hold the per-coordinate gap fixed, so $M = c\sqrt{\dd}$. The verified separation is
$M^2 Z = c^2 \dd Z$, and
\[
\dd Z = \frac{\chi^2_\kk}{(\chi^2_\kk + \chi^2_{\dd-\kk})/\dd} \;\Longrightarrow\; \chi^2_\kk
\]
as $\dd \to \infty$ with $\kk$ fixed, because the denominator converges to $1$ almost surely by the
law of large numbers while the numerator does not depend on $\dd$. By the continuous mapping theorem
and boundedness of $f$, expected acceptance converges to $\E[2 - 2\Phi(c\sqrt{\chi^2_\kk}/2)]$, which
contains no $\dd$. $\square$

So "acceptance is a function of $\kk$ alone" is a limit rather than a coincidence, with an $O(1/\dd)$
finite-size correction. Evaluating the exact Beta expression at $c = 0.578$ over
$\dd \in \{16, 64, 128, 512, 4096\}$ gives a coefficient of variation across dimensions of $0.15\%$ at
$\kk{=}1$, $0.47\%$ at $\kk{=}4$ and $1.7\%$ at $\kk{=}32$. The trained models show $1.3\%$--$12.3\%$
over a $32\times$ range of $\dd$, larger than the analytic prediction, which is what heteroscedastic
mixtures and finite evaluation samples should add.

\section{Audit of a Prior Continuous Speculative-Decoding Criterion}
\label{app:cspd}

This appendix is a code-level audit of one factor in the acceptance ratio of
\citet{wang2025continuous}, and we state its status carefully. We derive the trajectory ratio the
method's own construction implies, compare it against the released implementation, and bound the
size of the discrepancy. What we have \emph{not} done is reproduce that work's reported acceptance
rates or speedups end to end, or run a fidelity audit of its generated images, so we report an
inconsistency between a derivation and an implementation rather than a verdict on the method's
published results. We measure the relevant quantity on a real MAR-B/MAR-H pair below, and note that
the orientation in question is independently checkable from the public source.

\subsection{The trajectory ratio}

The method of \citet{wang2025continuous} verifies a token produced by a diffusion head. Crucially it
draws the diffusion noise once and shares it between the draft and target denoising trajectories, so
the two trajectories are coupled and the verified object is not the final sample but the entire
trajectory $x_T \to x_{T-1} \to \cdots \to x_0$. The correct density ratio is therefore
\[
\frac{p(x_{0:T})}{q(x_{0:T})} = \prod_{t=1}^{T} \frac{p(x_{t-1} \mid x_t)}{q(x_{t-1} \mid x_t)} ,
\]
since the shared $x_T$ contributes a factor of one. Each transition is Gaussian,
$p(x_{t-1}|x_t) = \mathcal{N}(\mu_{p,t}, \Sigma_{p,t})$ with diagonal $\Sigma$, so each term
contributes both an exponential and a normalising factor:
\[
\frac{p(x_{t-1}|x_t)}{q(x_{t-1}|x_t)}
= \frac{|\Sigma_{q,t}|^{1/2}}{|\Sigma_{p,t}|^{1/2}}
  \exp\!\Bigl(-\tfrac{1}{2}\|x_{t-1}-\mu_{p,t}\|^2_{\Sigma_{p,t}^{-1}}
              +\tfrac{1}{2}\|x_{t-1}-\mu_{q,t}\|^2_{\Sigma_{q,t}^{-1}}\Bigr).
\]
Collecting the normalisers over the trajectory gives an accumulated factor
\begin{equation}
\text{correct} \;=\; \prod_{t=1}^{T} \prod_{i=1}^{\dd} \frac{\sigma_{q,t,i}}{\sigma_{p,t,i}}
\;=\; \exp\Bigl(\textstyle\sum_{t,i} \log \sigma_{q,t,i} - \sum_{t,i} \log \sigma_{p,t,i}\Bigr).
\label{eq:correct-factor}
\end{equation}
Note the orientation: \emph{draft over target}, i.e.\ $q$ in the numerator, because the normaliser of
a Gaussian carries $\sigma^{-1}$ and the ratio inverts it.

\subsection{What the implementation computes}

The released implementation accumulates a per-model running sum of log standard deviations across the
sampler's steps and multiplies the ratio by
\begin{equation}
\text{shipped} \;=\; \exp\Bigl(\textstyle\sum_{t,i}\log\sigma_{p,t,i} - \sum_{t,i}\log\sigma_{q,t,i}\Bigr)
\;=\; \prod_{t,i} \frac{\sigma_{p,t,i}}{\sigma_{q,t,i}} .
\label{eq:shipped-factor}
\end{equation}
Two points of bookkeeping are worth stating because they are where an audit can go wrong. The
diffusion code returns $\tfrac{1}{2}\log \mathrm{Var}$ under the key \texttt{log\_variance}, so the
accumulator holds $\sum \log \sigma$ and not $\sum\log\sigma^2$; and the accumulation is over the
sampler's full schedule, so $T$ is the number of denoising steps, not one.

Comparing Eq.~\ref{eq:correct-factor} and Eq.~\ref{eq:shipped-factor}, the shipped factor is the
\emph{reciprocal} of the required one. It is not a spurious term: an independent derivation
\citep{zhang2025fastardiff} presents the same normaliser as part of the correct criterion, and our
analysis agrees that it belongs. It is applied upside down.

\subsection{Magnitude of the error}

Since $\text{shipped} = \text{correct}^{-1}$, using the shipped factor rather than omitting it moves
the log-ratio by $2\log(\text{correct})$: the error exponent is \emph{twice} the factor's own
exponent, and the ratio is wrong by $\text{correct}^2$. Bounding the accumulated variance mismatch
over the sampler's schedule and channel count gives:

\begin{center}
\small
\begin{tabular}{lr}
\toprule
quantity & value \\
\midrule
worst-case exponent of the factor (16 channels, full schedule) & $80.8$ \\
worst-case error exponent in the ratio ($2\times$) & $\mathbf{161.5}$ \\
exponent at an assumed $5\%$ per-step variance gap & $4.04$ \\
error exponent at that gap ($2\times$) & $8.08$, i.e.\ $\approx 3{,}200\times$ \\
exponent sufficient merely to double the acceptance rate & $0.693$ \\
\bottomrule
\end{tabular}
\end{center}

An exponent of $0.693$ doubles acceptance; the plausible operating point is an exponent of $8$. The
criterion therefore saturates $\min\{1,\cdot\}$ for the overwhelming majority of tokens and accepts
almost everything. On synthetic Gaussians an exponent of $4.0$ alone moves acceptance from $0.32$ to
$0.91$ at temperature $1$. A method that accepts nearly everything will report a large speedup, and
without an output-fidelity audit that speedup is indistinguishable from a correct one.

\subsection{Why this is a sign slip and not a modelling choice}

Two features of the source identify the orientation as accidental. First, the normaliser for the
\emph{final} denoising step appears in the same expression written the right way round, as a ratio of
draft to target standard deviations; only the accumulated term over the earlier steps is inverted.
Second, the accumulation is written as target-minus-draft while every other scale-carrying term in the
same file is draft-over-target. A deliberate design choice would be unlikely to be locally
inconsistent with the expression it sits inside.

We note that our own earlier analysis of this criterion claimed the accumulated factor had no place in
the ratio at all. That claim was wrong, and reading \citet{zhang2025fastardiff} against the source is
what corrected it. The resulting claim is weaker, since the term belongs and is merely inverted, and far more
defensible, with an independent derivation supporting the form.

\subsection{What the criterion is actually applied to, on a real MAR pair}
\label{app:e06}

We ran the released runner once on a genuine MAR-B draft against a MAR-H target
($208$M and $942$M parameters, ImageNet $256$, $64$ autoregressive steps, guidance $4.0$) and
recorded every input its acceptance test receives: $1008$ token decisions at $\dd = 16$. The
distributions it compares turn out to be nearly degenerate. The final-step posterior standard
deviation is $0.0064$ for both models, the two agree on it to a median ratio of $1.000$, and the
proposed token is the draft's own posterior mean to within numerical precision, because the last
denoising step adds no noise. The gap between the two means is $0.107$ in median absolute value, so
in units of that standard deviation the models sit a median $16.5\sigma$ apart per coordinate,
giving a Mahalanobis separation of $M \approx 138$ and an exact acceptance rate that rounds to zero.

This is a fact about the final denoising transition, not about the two models: the predictive
distribution a verifier should be comparing is the trajectory marginal, and the last transition of a
diffusion sampler is a near-deterministic map with a vanishing scale. We report it here because it
explains the role of the sampling temperature. At the shipped value of $42$ the effective separation
falls by the same factor, to $M \approx 3.3$, which is the regime in which the test discriminates at
all. Any acceptance rate this criterion reports is therefore inseparable from that choice, which is
the strongest reason we treat the audit as an appendix rather than a headline: what is
demonstrably inconsistent is the orientation of one factor, whereas assessing the method's reported
speedups would require a fidelity audit of its outputs that we have not run. Measuring the
trajectory-marginal separation of MAR-B and MAR-H, the quantity Figure~\ref{fig:feasibility}(b)
would place on the map, requires drawing from each head repeatedly at shared contexts and remains
future work.

\section{Trajectory-Marginal Separation of MAR-B and MAR-H}
\label{app:e07}

MAR's per-token density is produced by a diffusion head, so there is no $(\mu, \sigma)$ to read off.
The only Gaussian available in closed form is the final denoising transition, and
Appendix~\ref{app:e06} shows why that is the wrong object: its standard deviation is $0.006$ and its
draw is deterministic, which puts the two models a nominal $M \approx 138$ apart and says nothing
about the pair. What a speculative verifier compares is the trajectory marginal, and the honest way
to reach a distribution with no closed form is to draw from it.

\paragraph{Protocol.} One MAR-H generation supplies the trajectory. At $12$ autoregressive steps
spread across the schedule, the \emph{same} partially generated token map is handed to both models
--- the shared accepted prefix a speculative decoder would have --- and each model's head is sampled
$64$ times at up to $24$ masked positions. Per-coordinate means and standard deviations follow, and
with them $\delta$, $\sigma$, $M$, and the acceptance Proposition~\ref{prop:conservation} predicts.
The trajectory is advanced with the target's own draw throughout, so it remains a valid MAR-H
generation and the conditioning states are ones the sampler really visits.

\begin{center}
\small
\begin{tabular}{lcc}
\toprule
& guidance $3.0$ & guidance $1.0$ \\
\midrule
positions measured                    & 188 & 64 \\
$\sigma$ (target / draft)             & 0.4007 / 0.3984 & 0.4823 / 0.4841 \\
per-coordinate $|\delta|/\sigma$, mean & 0.377 & 0.282 \\
per-coordinate $|\delta|/\sigma$, median & 0.280 & 0.221 \\
separation $M$, mean                  & 1.850 & 1.400 \\
separation $M$, median [p5, p95]      & 1.683 [0.81, 3.57] & 1.256 [0.69, 2.60] \\
$M/\sqrt{\dd}$                        & 0.463 & 0.350 \\
predicted acceptance (95\% CI)        & $\best{0.397}$ [0.369, 0.425] & 0.513 [0.473, 0.550] \\
\bottomrule
\end{tabular}
\end{center}

\paragraph{Reading it.} Three points. The shared-covariance assumption that
Proposition~\ref{prop:conservation} needs is well supported for this pair: the two heads agree on
scale to within $0.6\%$ at guidance $3.0$ and $0.4\%$ without guidance, which is the regime in which
the conservation identity is exact rather than approximate. The separation is modest --- $M$ near
$1.9$, not the $138$ the final-step Gaussians suggest --- so exact verification of a single KL16
token is a real possibility rather than a lost cause. And guidance matters: raising it from $1.0$ to
$3.0$ pushes the pair apart from $M = 1.40$ to $1.85$ and costs $0.12$ in predicted acceptance,
because guidance is applied to both models but sharpens the larger one more.

The number is a prediction from the measured $M$ under a Gaussian model, not an acceptance rate
observed by running a speculative sampler; \S\ref{sec:limits} says what that leaves open. It is
enough to locate the pair on Figure~\ref{fig:feasibility}(b), which is what it is for.

\section{Distillation: the Divergence, Not Only the Acceptance}
\label{app:e15}

E1.4 reports that distillation raises exact acceptance at zero fidelity cost. That is the
consequence; the conservation argument names a cause, and a reviewer is right to ask whether the
cause was measured or assumed. It could have been assumed: a draft that merely inflated its own
variance would also raise overlap, and the acceptance table would look identical. So we re-ran the
same trajectory --- same checkpoints, schedule, seeds --- recording $\KL(p\|q)$ and $M$ at every
stage instead of the fidelity metrics. $\KL$ has no closed form for mixture heads, so it is estimated
by Monte Carlo, drawing from the frozen target and scoring under both; because the target is frozen,
$H(p)$ is constant and every change is a change in the draft.

\begin{center}
\small
\begin{tabular}{lccc ccc}
\toprule
& \multicolumn{3}{c}{$\KL(p\|q)$ (nats)} & \multicolumn{3}{c}{separation $M$} \\
\cmidrule(lr){2-4}\cmidrule(lr){5-7}
$\dd$ & before & after & reduction & before & after & reduction \\
\midrule
4   & 0.514  & 0.322  & \cgood{18}{37.3\%} & 0.320 & 0.212 & 33.8\% \\
8   & 1.674  & 1.269  & \cgood{12}{24.2\%} & 0.609 & 0.471 & 22.7\% \\
16  & 4.731  & 3.740  & \cgood{10}{20.9\%} & 1.075 & 0.830 & 22.8\% \\
32  & 9.724  & 7.822  & \cgood{10}{19.6\%} & 1.482 & 1.236 & 16.6\% \\
64  & 19.281 & 14.490 & \cgood{12}{24.8\%} & 1.922 & 1.551 & 19.3\% \\
128 & 37.860 & 28.710 & \cgood{12}{24.2\%} & 2.750 & 2.144 & 22.0\% \\
\bottomrule
\end{tabular}
\end{center}

The causal chain is visible link by link: $\KL(p\|q)$ down, $M$ down, acceptance up, at every
dimension. Acceptance rises monotonically across all five stages everywhere; $\KL$ and $M$ fall
monotonically except for wiggles below $1\%$ in the final stage at $\dd \in \{32, 64, 128\}$, which
are within the Monte Carlo error. The acceptance rates re-measured here ($0.0033 \to 0.0189$ at
$\dd{=}128$) agree with E1.4's independently measured $0.0036 \to 0.0208$.

\paragraph{Why a quarter buys $5.72\times$.} The reductions are modest and remarkably flat in
dimension --- around $20$--$25\%$ of the divergence at every width --- while the acceptance gain they
purchase grows from $1.12\times$ to $5.72\times$. That is the Gaussian tail doing the work, and it is
the clearest evidence we have for the mechanism of \S\ref{sec:moving}: the same proportional
reduction in the pair's separation is worth exponentially more where the tail is steeper.

\paragraph{One caveat on $M$.} For mixture heads $M$ is a summary of the first two moments, and we
compute it with the mixture's own total variance, which folds the spread \emph{between} components
into the scale. That makes $M$ smaller than the value governing acceptance, and correspondingly the
Gaussian formula $2 - 2\Phi(M/2)$ evaluated at it overpredicts: $0.372$ against a measured $0.0189$
at $\dd{=}128$. The absolute calibration is therefore not usable for mixtures, consistent with the
caveat in Appendix~\ref{app:proofs}, and we read $M$ here only as a relative indicator, before
against after. $\KL(p\|q)$ carries the argument because it needs no Gaussian assumption at all.

\section{Experimental Setup}
\label{app:setup}

\paragraph{Data.} ETTh1, the seven-channel hourly electricity-transformer series
\citep{zhou2021informer}. We split chronologically into $70/10/20$ train/validation/test fractions,
fit the standardiser on the training split only, and cut sliding windows of length $96$ \emph{after}
splitting, so that no window straddles a boundary and no test window shares a timestamp with a
training one. Test tokens number $3{,}389$ windows of $12$ tokens each.

\paragraph{Tokenizer.} A patch VAE \citep{kingma2014auto}, the one-dimensional analogue of the KL16
image tokenizer: a strided convolutional encoder maps a $(7, 96)$ window to $12$ tokens of dimension
$\dd$, and a mirrored decoder maps them back. Patch size $8$, hidden width $128$, three residual
blocks, KL weight $10^{-6}$. The posterior log-variance is clamped, because early in training the
encoder drives it towards $-60$ on easy patches and collapses the latent. Trained for $2{,}000$ steps
per dimension. Crucially, $\dd$ is a constructor argument and the tokenizer is retrained at every
$\dd$ in the sweep, which is what makes the acceptance-versus-$\dd$ curve possible at all.

\paragraph{Autoregressive models.} A masked autoregressive transformer over continuous tokens, the
one-dimensional analogue of MAR \citep{li2024autoregressive}. Masked positions are predicted from
unmasked ones, and the per-token distribution is a diagonal Gaussian mixture rather than a softmax
over a vocabulary. The mask ratio is drawn per sample from a truncated cosine distribution over
$[0.7, 1.0]$ of tokens, matching MAR: training predominantly on heavily masked inputs is what makes
the model usable at the early generation steps where almost nothing is known.

\begin{center}
\small
\begin{tabular}{lccccc}
\toprule
model & width & depth & heads & mixture components & role \\
\midrule
target & 256 & 6 & 8 & 8 & the expensive model being accelerated \\
draft  & 128 & 4 & 4 & 4 & the cheap proposal model \\
\bottomrule
\end{tabular}
\end{center}

Both are trained independently on the same data for $4{,}000$ steps with AdamW
\citep{loshchilov2019decoupled}, batch size $64$, and identical schedules. Training them independently
is deliberate: it is the configuration a practitioner reaches for when adding speculative decoding to
an existing model, and \S\ref{sec:moving} shows what changes when they are coupled.

\paragraph{Measurement.} All acceptance and fidelity numbers are computed on held-out test data
through a single shared estimator, so that no two experiments in the ladder can drift apart in how
they define acceptance. Approximately $25{,}000$ acceptance decisions contribute to each E1.1 point
and $\sim$$8{,}000$ to each E1.2b point. Random number generation is forked and re-seeded per
measurement: an earlier version seeded only the mask, and because \texttt{MixtureSameFamily.sample}
accepts no generator the token draw silently followed ambient global state, making measurements
irreproducible at the third decimal place.

\paragraph{Compute.} Every experiment in this paper runs on CPU. The models are small by design, with the
target under a million parameters, because the claims concern the token dimension and the
draft/target divergence, neither of which requires scale to exhibit. The full ladder reruns in under
an hour on eight threads.

\section{Convergence Controls for the Dimensional Collapse}
\label{app:convergence}

The obvious objection to Table~\ref{tab:e11} is that a fixed step budget leaves high-dimensional
models less converged, so the collapse could be a training artefact rather than a property of the
dimension. We address it twice.

\paragraph{Convergence traces.} Checkpointing throughout training at $\dd \in \{16, 128\}$ and
measuring acceptance between every pair of draft and target checkpoints shows both dimensions plateau
well inside the $4{,}000$-step budget: $\dd{=}16$ by roughly $750$ steps and $\dd{=}128$ by roughly
$500$. Extending to $16{,}000$ steps does not rescue acceptance; it \emph{falls} further, because
validation loss turns over early and the two models proceed to overfit apart rather than together.
More training makes speculative acceptance worse, not better, which is the opposite of what the
confound predicts.

\paragraph{Matched-$\delta$ comparison.} The sharper control dispenses with step counts entirely.
Across the checkpoint grid we find pairs at $\dd{=}16$ and $\dd{=}128$ whose measured per-coordinate
gap $\delta$ is equal, and compare their acceptance rates directly. At identical per-coordinate
agreement the $\dd{=}128$ models accept roughly $60\times$ less often: a median of $62\times$ over
$15$ matched pairs at the $4{,}000$-step budget and $64\times$ over $7$ pairs at $16{,}000$. Since
$\delta$ is held fixed by construction, the residual factor is attributable to dimension alone.

\section{Full Results}
\label{app:results}

\subsection{Random projections: the null control and the frontier}

Acceptance under a uniformly random $\kk$-dimensional statistic, averaged over five independent draws
of $A$ per cell, as a function of the verified dimension $\kk$ across all six token dimensions:

\begin{center}
\small
\begin{tabular}{lcccccccc}
\toprule
$\kk$ & 1 & 2 & 4 & 8 & 16 & 32 & 64 & 128 \\
\midrule
mean acceptance over $\dd$ & 0.8786 & 0.7846 & 0.6244 & 0.4102 & 0.2093 & 0.0774 & 0.0205 & 0.0035 \\
coefficient of variation & 1.3\% & 2.3\% & 2.5\% & 6.1\% & 12.3\% & 4.3\% & 13.5\% & n/a \\
number of $\dd$ contributing & 6 & 6 & 6 & 5 & 4 & 3 & 2 & 1 \\
\bottomrule
\end{tabular}
\end{center}

Acceptance is a function of $\kk$ essentially alone: at $\kk = 4$ the rate is $0.624$ whether the
ambient dimension is $4$ or $128$. As a check on the implementation, $\kk = \dd$ reproduces raw
verification to four decimal places at all six dimensions.

The acceptance/fidelity frontier for random projections, which is the baseline all structured rules
are scored against:

\begin{center}
\small
\begin{tabular}{lcccccccc}
\toprule
& \multicolumn{5}{c}{$\dd = 16$} & \multicolumn{3}{c}{$\dd = 64$ (abridged)} \\
\cmidrule(lr){2-6}\cmidrule(lr){7-9}
$\kk$ & 16 (exact) & 8 & 4 & 2 & 1 & 64 (exact) & 8 & 1 \\
\midrule
acceptance & 0.1613 & 0.3531 & 0.5960 & 0.7419 & 0.8475 & 0.0138 & 0.4094 & 0.8675 \\
excess NLL & $-0.003$ & 1.423 & 3.802 & 5.352 & 6.397 & $+0.125$ & 11.300 & 27.609 \\
\bottomrule
\end{tabular}
\end{center}

The exact anchors read $-0.1109 \pm 0.0912$ at $\dd{=}64$ and $+0.1883 \pm 0.1896$ at $\dd{=}128$, both
within $1.2$ standard errors of zero and consistent with the losslessness guarantee of
Appendix~\ref{app:exactness}.

\subsection{Five subspace rules at matched $\kk$}

The complete E1.2b measurement. Acceptance, excess NLL in nats per token, and signal-domain energy
distance, for every rule and every $\kk$. Recall from Proposition~\ref{prop:rowspace} that each row is
a choice of subspace; the rules differ enormously at matched $\kk$ and yet, when replotted against
acceptance, trace one curve (Table~\ref{tab:collapse}).

\begin{center}
\small
\begin{tabular}{llccc|llccc}
\toprule
\multicolumn{5}{c|}{$\dd = 64$} & \multicolumn{5}{c}{$\dd = 128$} \\
$\kk$ & rule & acc. & nats & energy & $\kk$ & rule & acc. & nats & energy \\
\midrule
64 & exact         & 0.018 & $-0.11$ & $-0.14$ & 128 & exact         & 0.004 & $+0.19$ & $+0.15$ \\
\midrule
1 & random        & 0.868 & $+25.04$ & $+1.05$ & 1 & random        & 0.886 & $+46.83$ & $+1.16$ \\
1 & pca           & 0.614 & $+17.15$ & $+0.48$ & 1 & pca           & 0.525 & $+26.44$ & $+0.46$ \\
1 & signal-top    & 0.668 & $+19.75$ & $+0.76$ & 1 & signal-top    & 0.592 & $+31.27$ & $+0.60$ \\
1 & dct-low       & 0.666 & $+19.69$ & $+0.77$ & 1 & dct-low       & 0.592 & $+31.00$ & $+0.59$ \\
1 & signal-bottom & 0.918 & $+26.45$ & $+1.34$ & 1 & signal-bottom & 0.916 & $+48.40$ & $+1.44$ \\
\midrule
2 & random        & 0.799 & $+22.78$ & $+0.86$ & 2 & random        & 0.794 & $+41.49$ & $+0.88$ \\
2 & pca           & 0.359 & $+8.76$ & $+0.26$ & 2 & pca           & 0.299 & $+13.72$ & $+0.25$ \\
2 & signal-top    & 0.441 & $+12.05$ & $+0.16$ & 2 & signal-top    & 0.267 & $+11.91$ & $+0.23$ \\
2 & dct-low       & 0.438 & $+11.96$ & $+0.11$ & 2 & dct-low       & 0.268 & $+11.96$ & $+0.25$ \\
2 & signal-bottom & 0.863 & $+24.50$ & $+1.34$ & 2 & signal-bottom & 0.865 & $+44.99$ & $+1.13$ \\
\midrule
4 & random        & 0.652 & $+17.74$ & $+0.56$ & 4 & random        & 0.652 & $+32.78$ & $+0.65$ \\
4 & pca           & 0.144 & $+2.78$ & $+0.01$ & 4 & pca           & 0.084 & $+2.80$ & $+0.20$ \\
4 & signal-top    & 0.182 & $+3.70$ & $+0.02$ & 4 & signal-top    & 0.114 & $+4.48$ & $+0.15$ \\
4 & dct-low       & 0.181 & $+3.67$ & $+0.01$ & 4 & dct-low       & 0.113 & $+4.41$ & $+0.15$ \\
4 & signal-bottom & 0.791 & $+21.40$ & $+0.99$ & 4 & signal-bottom & 0.796 & $+40.58$ & $+0.83$ \\
\midrule
8 & random        & 0.419 & $+10.24$ & $+0.12$ & 8 & random        & 0.432 & $+20.19$ & $+0.26$ \\
8 & pca           & 0.075 & $+1.05$ & $-0.15$ & 8 & pca           & 0.034 & $+1.11$ & $+0.17$ \\
8 & signal-top    & 0.100 & $+1.66$ & $-0.10$ & 8 & signal-top    & 0.042 & $+1.33$ & $+0.17$ \\
8 & dct-low       & 0.098 & $+1.63$ & $-0.11$ & 8 & dct-low       & 0.041 & $+1.33$ & $+0.16$ \\
8 & signal-bottom & 0.683 & $+17.67$ & $+0.70$ & 8 & signal-bottom & 0.685 & $+33.62$ & $+0.75$ \\
\midrule
16 & random        & 0.214 & $+4.21$ & $-0.06$ & 16 & random        & 0.268 & $+11.25$ & $+0.08$ \\
16 & pca           & 0.047 & $+0.48$ & $-0.12$ & 16 & pca           & 0.018 & $+0.56$ & $+0.14$ \\
16 & signal-top    & 0.056 & $+0.59$ & $-0.14$ & 16 & signal-top    & 0.024 & $+0.77$ & $+0.16$ \\
16 & dct-low       & 0.059 & $+0.69$ & $-0.14$ & 16 & dct-low       & 0.025 & $+0.77$ & $+0.16$ \\
16 & signal-bottom & 0.520 & $+12.35$ & $+0.20$ & 16 & signal-bottom & 0.550 & $+25.43$ & $+0.43$ \\
\midrule
32 & random        & 0.079 & $+0.91$ & $-0.11$ & 32 & random        & 0.090 & $+2.96$ & $+0.13$ \\
32 & pca           & 0.033 & $+0.14$ & $-0.09$ & 32 & pca           & 0.013 & $+0.39$ & $+0.16$ \\
32 & signal-top    & 0.037 & $+0.21$ & $-0.14$ & 32 & signal-top    & 0.016 & $+0.45$ & $+0.14$ \\
32 & dct-low       & 0.038 & $+0.27$ & $-0.08$ & 32 & dct-low       & 0.017 & $+0.45$ & $+0.17$ \\
32 & signal-bottom & 0.314 & $+6.56$ & $-0.02$ & 32 & signal-bottom & 0.345 & $+14.80$ & $+0.04$ \\
\bottomrule
\end{tabular}
\end{center}

Reading the table by $\kk$ suggests five very different methods: at $\dd{=}64, \kk{=}4$, acceptance
ranges from $0.14$ to $0.79$, a factor of $5.6$. Reading it by acceptance shows one curve. Note also
the near-identity of \textbf{signal-top} and \textbf{dct-low} at every cell; the low-frequency band of
the decoded patch and the decoder's overall sensitivity select nearly the same subspace, which is a
statement about this decoder rather than about acceptance rules in general.

\subsection{The oracle-disagreement subspace}
\label{app:oracle}

The five rules of E1.2b all guess where the draft and target differ. This control removes the guess.
For each dimension we collect the per-decision mean gap $\delta = \mu_p - \mu_q$ on the same masked
decisions the fidelity measurement scores, form the uncentred second moment
$\E[\delta\delta^\top]$, and verify its top-$\kk$ eigenvectors. By construction that is the rank-$\kk$
subspace capturing the largest share of $\E\|\delta\|^2$, so it upper-bounds what \emph{any}
rank-$\kk$ linear statistic can capture of the disagreement. The second moment is uncentred because
what the test sees at each decision is the gap itself, not its deviation from an average gap.

The rule is an oracle in the strict sense: building it requires the target's predictions on the
decisions being verified, which is exactly what a verifier does not have before it runs the target.
No one can deploy this. Its value is that it closes the escape route --- if the frontier were an
artefact of five poorly chosen heuristics, the rule that cannot be poorly chosen should break it.

\begin{center}
\small
\begin{tabular}{rrccrr}
\toprule
$\dd$ & $\kk$ & acceptance & excess NLL & $\E\|\delta\|^2$ captured & residual vs.\ random \\
\midrule
4   & 1 & 0.882 & $+0.55$  & 39.9\% & $-0.002$ \\
4   & 2 & 0.773 & $+0.36$  & 72.6\% & --- \\
\midrule
8   & 1 & 0.826 & $+1.97$  & 35.1\% & $+0.082$ \\
8   & 2 & 0.706 & $+1.40$  & 62.0\% & $+0.050$ \\
8   & 4 & 0.556 & $+0.79$  & 87.8\% & --- \\
\midrule
16  & 1 & 0.770 & $+5.38$  & 38.7\% & $+0.107$ \\
16  & 2 & 0.548 & $+3.29$  & 58.6\% & $+0.181$ \\
16  & 4 & 0.371 & $+1.61$  & 86.4\% & --- \\
\midrule
32  & 1 & 0.720 & $+9.44$  & 32.1\% & $+0.393$ \\
32  & 2 & 0.463 & $+5.54$  & 57.7\% & $+0.535$ \\
32  & 4 & 0.264 & $+2.44$  & 81.9\% & $+0.300$ \\
\midrule
64  & 1 & 0.654 & $+19.21$ & 28.1\% & $+1.395$ \\
64  & 2 & 0.379 & $+10.22$ & 49.8\% & $+1.158$ \\
64  & 4 & 0.152 & $+2.85$  & 77.6\% & $+0.147$ \\
\midrule
128 & 1 & 0.534 & $+27.42$ & 36.2\% & $+1.362$ \\
128 & 2 & 0.239 & $+10.69$ & 55.9\% & $+0.815$ \\
128 & 4 & 0.084 & $+2.96$  & 78.0\% & --- \\
\bottomrule
\end{tabular}
\end{center}

Residuals are against E1.2b's random-projection curve interpolated to the same acceptance, with no
extrapolation; a dash marks a point whose acceptance falls outside the range that curve was measured
over, and larger $\kk$ (omitted for space) is always in that category. A \emph{positive} residual is
the oracle costing more than a random subspace at the same acceptance.

Two things stand out. The disagreement is concentrated --- one direction carries $28$--$40\%$ of
$\E\|\delta\|^2$ and four carry $78$--$91\%$ --- so this is not a case of the rule having nothing to
find. And it still does not pay: every residual outside $\dd{=}4$ is positive, meaning the oracle is
if anything \emph{worse} than a structureless projection at matched acceptance, and the largest
departure anywhere is $5.3\%$ of the cost range ($\dd{=}64$), inside the $4.0$--$15.9\%$ envelope the
five heuristics occupy in Table~\ref{tab:collapse}. The concentration of $\delta$ buys a favourable
position \emph{along} the curve, not a position off it, which is what
Corollary~\ref{cor:frontier} predicts: capturing more of $\delta$ lowers the leak and lowers
acceptance by exactly the amount that keeps the product on the frontier.

\subsection{Uncertainty in the collapse statistic}
\label{app:bootstrap}

The collapse statistic quoted in Table~\ref{tab:collapse} is a ratio of two numbers each estimated
with error, so it deserves an interval. The bootstrap is parametric, because
\texttt{results.json} stores per-cell means and standard errors rather than raw per-decision values
and there is therefore nothing to resample nonparametrically. Each of $10^4$ replicates perturbs
every cost by its reported standard error and every acceptance by its binomial standard error
$\sqrt{a(1-a)/n}$, then recomputes the statistic from scratch --- including re-interpolating the
random-projection baseline at the perturbed acceptances, so that uncertainty in the curve the
structured rules are \emph{scored against} is propagated and not just uncertainty in the scores
(\texttt{scripts/bootstrap\_frontier.py}).

\begin{center}
\small
\begin{tabular}{lcccc}
\toprule
& \multicolumn{2}{c}{excess NLL} & \multicolumn{2}{c}{energy distance} \\
\cmidrule(lr){2-3}\cmidrule(lr){4-5}
$\dd$ & max (\% of range) & median & max (\% of range) & median \\
\midrule
4   & 15.9 \, [7.4, 30.4] & 11.1 \, [3.0, 20.7] & 36.1 \, [22.5, 87.1] & 32.7 \, [8.3, 55.0] \\
8   & 7.2 \, [4.6, 11.2]  & 1.1 \, [0.9, 3.9]   & 36.8 \, [29.1, 60.8] & 23.2 \, [9.9, 31.8] \\
16  & 6.2 \, [4.7, 8.6]   & 0.9 \, [0.8, 2.7]   & 21.9 \, [18.2, 38.2] & 8.3 \, [5.7, 17.3] \\
32  & 8.4 \, [6.7, 10.2]  & 2.9 \, [2.1, 4.0]   & 16.1 \, [15.0, 35.8] & 7.5 \, [4.2, 15.2] \\
64  & 5.6 \, [5.0, 7.5]   & 2.3 \, [1.7, 3.2]   & 20.6 \, [14.4, 35.6] & 5.5 \, [4.0, 12.5] \\
128 & 4.0 \, [3.3, 5.6]   & 1.6 \, [1.2, 2.5]   & 11.6 \, [12.7, 34.8] & 3.5 \, [3.5, 13.5] \\
\bottomrule
\end{tabular}
\end{center}

The excess-NLL intervals are tight enough to carry the argument. At $\dd{=}128$ the maximum departure
from a single curve is $4.0\%$ of the range with a $95\%$ interval of $3.3$--$5.6\%$, and that
interval is disjoint from the one at $\dd{=}4$ ($7.4$--$30.4\%$), so the tightening of the collapse
with dimension is resolved rather than an artefact of noise. The median departure is smaller
still, $1.6\%$ at $\dd{=}128$: the maximum is driven by one or two cells, and typical behaviour is
closer to the curve than the headline number suggests.

\paragraph{Two cautions, both about the maximum.} Bootstrapping an extremum is upward-biased, because
the maximum of a set of noisy estimates exceeds the maximum of their means in almost every replicate.
The intervals on the max columns should therefore be read as conservative. Where the bias is severe
enough that the point estimate falls below its own interval --- energy distance at $\dd{=}128$, point
$11.6$ against interval $[12.7, 34.8]$ --- the explanation is that bias and not a disagreement
between the two computations. This is also why the median columns are reported: they are the stable
statistic, and they tell the same story with room to spare. The energy-distance intervals are wide
throughout, which is the same higher variance that makes us treat excess NLL as the primary metric
(Appendix~\ref{app:metric}) rather than a defect specific to this analysis.

\subsection{Distillation}

Full stage-wise results for the alignment experiment of \S\ref{sec:moving}. Excess NLL remains zero
within error at every stage, confirming that the acceptance gains are free rather than traded:

\begin{center}
\small
\begin{tabular}{lccccccc}
\toprule
$\dd$ & 0 steps & 2{,}000 & 4{,}000 & 6{,}000 & 8{,}000 & 10{,}000 & total gain \\
\midrule
4 & 0.6453 & 0.6860 & 0.6977 & 0.7131 & 0.7078 & 0.7198 & \cgood{12}{$\best{1.12\times}$} \\
8 & 0.4018 & 0.4383 & 0.4542 & 0.4600 & 0.4680 & 0.4728 & \cgood{15}{$\best{1.18\times}$} \\
16 & 0.1779 & 0.2095 & 0.2244 & 0.2311 & 0.2431 & 0.2507 & \cgood{20}{$\best{1.41\times}$} \\
32 & 0.0722 & 0.0891 & 0.0960 & 0.1040 & 0.1017 & 0.1084 & \cgood{23}{$\best{1.50\times}$} \\
64 & 0.0176 & 0.0251 & 0.0328 & 0.0388 & 0.0413 & 0.0459 & \cgood{32}{$\best{2.61\times}$} \\
128 & 0.0036 & 0.0073 & 0.0120 & 0.0142 & 0.0210 & 0.0208 & \cgood{44}{$\best{5.72\times}$} \\
\bottomrule
\end{tabular}
\end{center}

\section{The Signal-Domain Metric, and the One That Failed}
\label{app:metric}

Reporting a fidelity cost for a relaxed acceptance rule requires a metric that a bad rule cannot win.
Our first choice failed this requirement, and we record the failure because the failure mode is
general and quiet.

\paragraph{Excess squared error rewards mode collapse.} The natural signal-domain metric is: decode
the emitted tokens, measure squared error against the real window, and subtract the error a genuine
target draw would have incurred. Under it we measured, at $\dd{=}16$, $\kk{=}4$:

\begin{center}
\small
\begin{tabular}{lccc}
\toprule
rule & acceptance & excess NLL & excess MSE \\
\midrule
exact         & 0.144 & $-0.087$ & $-0.0068$ \\
signal-top    & 0.383 & $+1.880$ & $-0.0152$ \\
signal-bottom & 0.712 & $+5.122$ & $\mathbf{-0.0636}$ \\
\bottomrule
\end{tabular}
\end{center}

Every rule scores \emph{negative} excess distortion, apparently improving on genuine target
samples, and the anti-control, which by construction verifies the least important directions, scores
best of all by a factor of four. The explanation is that the draft is the smaller model, its
predictive distribution is tighter, and its samples therefore sit closer to the conditional mean.
Squared error against a single ground-truth reference is minimised by the conditional mean, not by a
correct sample from the predictive distribution, so any rule that emits more draft samples scores
better. A degenerate rule emitting the conditional mean every time would win outright.

\paragraph{Energy distance.} We replace it with the two-sample energy statistic
\citep{szekely2013energy}, computed on decoded waveforms under the same conditioning:
\[
\mathcal{E} = 2\,\E\|X - Y\| - \E\|X - X'\| - \E\|Y - Y'\|,
\]
with $X, X'$ two independent emissions of the rule and $Y, Y'$ two independent target draws. It is
non-negative, and zero exactly when the two distributions agree. The $-\E\|X - X'\|$ term rewards
spread, so shrinking the predictive variance \emph{increases} the statistic rather than decreasing it.
Re-measuring the same three rules:

\begin{center}
\small
\begin{tabular}{lccc}
\toprule
rule & acceptance & excess NLL & energy distance \\
\midrule
exact         & 0.135 & $-0.026$ & $+0.060 \pm 0.250$ \\
random        & 0.503 & $+3.229$ & $+0.337$ \\
signal-top    & 0.353 & $+1.953$ & $+0.069$ \\
signal-bottom & 0.724 & $+5.409$ & $\mathbf{+1.313}$ \\
\bottomrule
\end{tabular}
\end{center}

The ordering is now the one the construction demands: the anti-control is decisively worst, and exact
verification anchors at zero. We verify the estimator's properties directly in unit tests (zero for
matching distributions, positive under a mean shift, and positive under variance shrinkage) rather
than relying on the pipeline to exhibit them.

\paragraph{Both metrics are reported everywhere.} Excess NLL asks whether we are still sampling from
the model; the energy distance asks whether the reconstructed waveform still has the right
distribution. They can disagree, and a method reporting only one is easy to fool. The conclusions of
\S\ref{sec:conservation} hold under both.

\section{What Would Falsify These Claims}
\label{app:falsifiers}

We state the falsifiers explicitly, since a conservation claim is a claim that something is
impossible and is therefore only as strong as the effort made to break it.

\begin{itemize}
\itemsep2pt
\item \textbf{The collapse is a training artefact.} Falsified if models trained to convergence at
matched validation loss show acceptance independent of $\dd$. Tested in
Appendix~\ref{app:convergence}: matched-$\delta$ pairs still show a $60\times$ penalty, and further
training makes acceptance worse.

\item \textbf{Some subspace beats the frontier.} Falsified by any rule sitting materially below the
random-projection curve at matched acceptance. Five rules were tried, including one designed with full
knowledge of the decoder and one designed to be as bad as possible; the largest departure in either
direction is $4.0\%$ of the range, and the sign of the departure is not stable across $\kk$. A
nonlinear statistic is not covered by Proposition~\ref{prop:rowspace} and remains open, though it
forfeits the exactness of the projected test.

\item \textbf{The fidelity metric is measuring the harness.} Falsified if exact verification scores as
lossy. It does not: excess NLL is $-0.008 \pm 0.080$ and energy distance $-0.026$ at $\dd{=}64$, with
the residual sampler converging on every rejection. This anchor is checked in every experiment.

\item \textbf{The fidelity metric is gameable.} This one \emph{fired} during the work, against squared
error, and Appendix~\ref{app:metric} reports it. The replacement is chosen specifically to be immune
to the failure that occurred, and that immunity is unit-tested rather than assumed.

\item \textbf{Distillation gains are traded, not free.} Falsified if excess NLL rises as acceptance
does. It does not; the rule remains exact and Appendix~\ref{app:exactness} guarantees losslessness
analytically, with the measurement confirming it at every stage.
\end{itemize}

\section{When Can Speculative Sampling Work in a Continuous Token Space?}
\label{app:practical}

Proposition~\ref{prop:conservation} is usable as a design rule, not only as an obstruction. Because
acceptance depends on the pair through the single scalar $M = \|\delta\|/\sigma$, a practitioner can
predict the achievable speedup before implementing anything.

\paragraph{The budget.} Speculative sampling with a draft costing a fraction $c$ of the target's
forward pass yields an expected speedup of roughly $1/(c + 1 - \bar{A})$ for single-token
verification, so a rate of $\bar{A} = 0.5$ is worth having and $\bar{A} = 0.02$ is worth nothing.
Inverting $\bar{A} = 2 - 2\Phi(M/2)$ at $\bar A = 0.5$ requires $\Phi(M/2) = 0.75$, so useful
acceptance requires
\[
M = \frac{\|\delta\|}{\sigma} \;=\; \frac{\delta_{\text{coord}}\sqrt{\dd}}{\sigma}
\;\lesssim\; 1.35 .
\]
At $\dd = 16$ this permits a per-coordinate relative disagreement of about $0.34\sigma$; at
$\dd = 256$ it permits $0.084\sigma$; at $\dd = 4096$, $0.021\sigma$. Figure~\ref{fig:feasibility}(b)
plots the same constraint at four acceptance levels. The requirement tightens as $1/\sqrt{\dd}$, and
it is a requirement on how closely two \emph{separately trained} networks agree pointwise in latent
space, a far stronger condition than matching in distribution, and one that nothing in ordinary
training encourages.

\paragraph{Reading our measurements against it.} Our models sit at
$\delta_{\text{coord}}/\sigma \approx 0.2$, which is a good draft by any ordinary standard: the two
networks agree to within a fifth of a standard deviation on every coordinate. That is sufficient at
$\dd = 4$ ($M = 0.93$, $\bar A = 0.64$) and hopeless at $\dd = 128$ ($M = 5.83$,
$\bar A = 0.0035$). The same drafter is excellent or useless depending only on the width of the
latent it operates in.

\paragraph{Consequences for practice.} Three follow directly.
\begin{itemize}
\itemsep2pt
\item \textbf{Narrow latents are worth engineering for.} A tokenizer emitting $\dd = 8$ admits
speculative decoding on drafts that $\dd = 64$ would reject outright. This is a reason to prefer
narrow latents that has nothing to do with reconstruction quality, and it is not visible from any
single-$\dd$ pipeline.
\item \textbf{Do not budget for a projected rule.} It will raise acceptance, and a random one raises it
$180\times$ at $\dd{=}128$, but Table~\ref{tab:collapse} shows the fidelity paid is the same
whichever subspace is chosen, so the choice is not a design lever and effort spent selecting it is
wasted. If a relaxation is acceptable at all, the cheapest correct one is a random projection at the
$\kk$ that meets the fidelity budget.
\item \textbf{Couple the pair during training.} \S\ref{sec:moving} shows post-hoc distillation moves
the frontier, and it is the weakest form of coupling available, since the draft chases a target it never
influenced. Architectures that already maintain a predictor and an EMA target encoder produce a
matched pair as a by-product of training, and our results suggest that is where the leverage is.
\end{itemize}

\paragraph{What we have not shown.} That coupled training closes the gap at large $\dd$. Our
distillation gains are real but small, and the $\sqrt{\dd}$ requirement above suggests they must
become harder, not easier, as $\dd$ grows: halving $M$ at $\dd = 128$ demands halving the
per-coordinate gap everywhere at once. Whether any training scheme achieves that is the open
question our results leave, and it is the one worth attacking next.

\section{Reproducibility}
\label{app:repro}

Every number in this paper is produced by a scripted experiment that records its own git commit,
Python and PyTorch versions, full argument list, and wall-clock time into a JSON artefact alongside
the measurements it reports. Acceptance and fidelity are computed by single shared estimators rather than
per-experiment reimplementations, so the rungs of the ladder cannot drift apart in what they measure.
The implementation carries $113$ tests, including the losslessness anchor of
Appendix~\ref{app:exactness}, the $\kk = \dd$ bijection identity, row-space invariance
(Proposition~\ref{prop:rowspace}), the conservation identity (Proposition~\ref{prop:conservation}),
the exactness of the DCT band decomposition via Parseval, and the variance-shrinkage property of the
energy distance. All experiments run on CPU in under an hour.
